\newpage
\section*{Rezumat}

Această lucrare de dizertație explorează posibilitatea dezvoltării unui sistem de neurofeedback multi-utilizator conceput pentru sălile de clasă, utilizând dispozitive Brain-Computer Interface (BCI) cu un singur electrod și cost redus --- tipul de hardware accesibil vizat de sistem, precum NeuroSky MindWave Mobile 2. Pentru prototipul în sine a fost utilizată o cască g.tec medical Unicorn Hybrid Black cu 8 canale disponibilă la universitatea noastră; deși este un dispozitiv mai scump, pentru metrica de atenție este procesat un singur canal, astfel că ambele sunt la fel de potrivite pentru sarcină. Sistemul vizează colectarea și procesarea datelor de atenție de la mai mulți studenți în timp real, oferind educatorilor o perspectivă valoroasă asupra concentrării și implicării studenților pe parcursul diferitelor segmente ale unui curs.

Arhitectura sistemului este construită în jurul unei rețele de căști BCI care transmit date prin protocolul Lab Streaming Layer (LSL) către un calculator operat de profesor --- un echipament pe care majoritatea personalului didactic îl are deja la dispoziție. Va fi implementată și o aplicație Python dedicată, cu scopul de a vizualiza dinamic nivelul de atenție al fiecărui student, permițând instructorului să identifice nu doar tipare de implicare, ci și de dezimplicare. Sistemul permite, de asemenea, structurarea cursurilor în secțiuni predefinite, corelând nivelurile de atenție cu segmente specifice ale materialului. Aceasta facilitează analiza post-curs pentru a determina care subiecte sau segmente au susținut cel mai mult (sau cel mai puțin) interesul studenților, sprijinind educatorii în efectuarea de ajustări bazate pe date în strategiile lor de predare.

Spre deosebire de majoritatea lucrărilor pe această temă, accentul acestei teze nu se pune pe analiza rezultatelor învățării sau pe compararea metodologiilor de predare, ci pe evaluarea fezabilității tehnice și practice a unui astfel de sistem într-un context educațional și a modului în care acesta poate influența procesul de învățare în ansamblu. Lucrarea investighează aspecte precum sincronizarea datelor, calitatea semnalului, scalabilitatea și experiența utilizatorului, oferind o bază pentru cercetări viitoare în educația bazată pe neurofeedback cu ajutorul acestui sistem. În final, acest proiect urmărește să contribuie la dezvoltarea unor instrumente educaționale inteligente care valorifică neurotehnologia pentru a îmbunătăți eficiența predării și gradul de implicare al studenților.
